% !TEX root = template.tex

% \AtBeginSection[]{
%     \begin{frame}{目次}
%         \tableofcontents[currentsection, hideallsubsections]
%     \end{frame}
% }

\def\SlideTitle{ランダムDFT拡散スペクトル有効周波数分割多重と期待値伝播法に基づく低計算量復調法}

\author[小松和暉]{小松和暉}
\title[\SlideTitle（6月RCS研究会@北海道大学）]{\SlideTitle}
\institute[豊橋技科大]{豊橋技術科学大学}
\date{\today}

\begin{frame}[t,plain,noframenumbering]{}
% ロゴを入れる
\begin{tikzpicture}[overlay]
\node at (\paperwidth-12ex,-1.2cm) {
\includegraphics[width=7ex]{logo/tut_logo.pdf}
};
\end{tikzpicture}
%

\vspace{2cm}
% \vfill

\begin{rightnote}
\begin{center}
{
\Large
ランダムDFT拡散スペクトル有効周波数分割多重と\\期待値伝播法に基づく低計算量復調法
}
\\
{
\small
Randomized DFT Spread Spectrally Efficient Frequency Division Multiplexing\\
and Low-Complexity Expectation-Propagation Detection
}
\end{center}
\end{rightnote}

% \vspace{0.5cm}
\vfill

\begin{center}
\begin{tabular}{rrl}
$\dagger$豊橋技術科学大学&〇&小松和暉 \\
$\dagger\dagger$東京大学&&小島駿 \\
$\dagger$&&上原秀幸 \\
\end{tabular}

% \vspace{0.5cm}
\vfill

{
電子情報通信学会　無線通信システム研究会\\
6月研究会@北海道大学
}
\end{center}
\end{frame}


% \framewithnote{フーリエ変換とその逆変換}
% {
% \headuline{CA}{\textbf{フーリエ変換}}

% \UnderlineB{信号$x(t)$}から\UnderlineA{周波数成分$X(f)$}を求める積分変換
% \begin{align*}
% \UnderlineA{$X(f)$} = \int_{-\infty}^{\infty} \UnderlineB{$x(t)$} \mathrm{e}^{-j2\pi ft} \mathrm{d}t
% \end{align*}

% \vspace{1em}
% \headuline{CA}{\textbf{逆フーリエ変換}}

% \UnderlineA{周波数成分$X(f)$}から\UnderlineB{信号$x(t)$}を再構成する積分変換
% \begin{align*}
%     \UnderlineB{$x(t)$} = \int_{-\infty}^{\infty} \UnderlineA{$X(f)$} \mathrm{e}^{j2\pi ft} \mathrm{d}f
% \end{align*}

% }
% {
% \rightnotebox{メモ1}{\\
% \UnderlineA{周波数成分$X(f)$}を信号の\textbf{スペクトル}と呼ぶ
% }
% \rightnotebox{注意1}{\\
% 順変換と逆変換で$\mathrm{e}$の肩に乗っている$j2\pi f t$の符号が異なることに注意
% }
% }


% \framewithnote{その2}{
% \begin{infobox}{あいうえお}
% \EmpA{かき}くけこ
% \end{infobox}
% \begin{alertbox}{さしすせそ}
% たちつてと
% \end{alertbox}

% % \headuline{\textcolor{CA}{みだし}\hfill}
% \headuline{CA}{\textbf{逆フーリエ変換}}


% \headuline{CB}{みだし}
% }{}





\begin{frame}[t]{スペクトル有効周波数分割多重：SEFDM}
\vspace{-0.5em}
\doublecolumn[0.48]{
\scalebox{0.9}{
\input{figs/sefdm_ofdm.tex}
}
}{
\small
\scalebox{0.9}{
\input{figs/sefdm_mod.tex}
}
}


\diffdoublecolumn{0.78}{0.27}{
\begin{itemize}
\item 直交周波数分割多重（OFDM）より\TextCA{狭い帯域幅で\\同一の伝送速度}で情報を伝送する技術
\item OFDMのように逆高速フーリエ変換（IFFT）を\\用いて\TextCA{低計算量で変調が可能}
\end{itemize}
}{
\rightnotebox{Note}{
\begin{align*}
&\text{圧縮率 } \alpha = N/K < 1 \\
&\text{サブキャリア数} N \\
&\text{FFTサイズ} K
\end{align*}
}
}
\end{frame}




\begin{frame}[t]{SEFDMにおけるキャリア間干渉（ICI）}
\vspace{-2em}

\doublecolumn[0.48]{
\begin{center}
\input{figs/sefdm_ici_ofdm.tex}\\
サブキャリアが直交しているため\\FFTで簡単に復調可能
\end{center}
}{
\begin{center}
\input{figs/sefdm_ici_sefdm.tex}\\
サブキャリアが直交しておらず\\
\TextCB{ICIによりFFTでは復調できない}
\end{center}
}


\vspace{1em}
\singlecolumn[0.06]{
\begin{infobox}{SEFDMの特徴（OFDMとの比較）}
\begin{itemize}
\item サブキャリアを密に配置するため\TextCA{必要な帯域幅が減らせる}
\item \TextCB{ICIにより復調が困難}
\end{itemize}
\end{infobox}
}



% \begin{itemize}
%     \item OFDMでは
%     \item SEFDMでは
% \end{itemize}
\end{frame}



% \begin{frame}[t]{OFDMとSEFDMの比較（まとめ）}
%     \vspace{-1em}
%     \begin{center}
%         \renewcommand\arraystretch{2}
%         \tabcolsep = 0.5cm
%         \begin{tabular}{cccc}
%             & 帯域幅 & 変調 & 復調 \\ \hline
%             OFDM & $BW$ & IFFT & FFT \\ \hline
%             SEFDM & \TextCA{$\alpha BW$}; $(\alpha < 1)$ & IFFT & \TextCB{困難} \\ \hline
%         \end{tabular}
%     \end{center}

%     \begin{infobox}{SEFDMの特徴}
%         \begin{itemize}
%             \item \TextCA{必要な帯域幅を減らしつつ同一の伝送速度を達成}
%             \item \TextCB{ICIにより復調が困難}
%         \end{itemize}
%     \end{infobox}
% \end{frame}


\begin{frame}[t]{既存のSEFDMの復調方法}
\vspace{-1.5em}
\begin{center}
\renewcommand\arraystretch{2}
\tabcolsep = 0.1cm
\begin{tabular}{cccc}
基礎技術 & 復調性能 & 計算量 & 問題点 \\ \hline
最尤推定 & ○○ & ×× & 指数計算量 \\ \hline
固有モード伝送 & ○ & ○ & 小さな特異値の問題 \\ \hline
確率推論 & ○ & ○ & 小さな$N$で性能劣化 \\ \hline
\end{tabular}
\end{center}

\vspace{0.2em}
\singlecolumn[0.11]{
\begin{infobox}{}
\begin{itemize}
\item 基本的にはMIMO伝送の復調方式と同等な\\アルゴリズムによって復調可能
\item それぞれ一長一短があるのが現状
\end{itemize}
\end{infobox}
}
\end{frame}



\begin{frame}[t]{大規模SEFDMにおける復調方法の比較}
\vspace{-1.5em}
\begin{center}
\renewcommand\arraystretch{2}
\tabcolsep = 0.1cm
\begin{tabular}{cccc}
基礎技術 & 復調性能 & 計算量 & 問題点 \\ \hline
最尤推定 & ○○ & ×× & \TextCB{指数計算量} \\ \hline
固有モード伝送 & ○ & ○ & \TextCB{小さな特異値の問題} \\ \hline
\TextCA{確率推論}　& ○ & ○ & \textcolor{black!20}{小さな$N$で性能劣化} \\ \hline
\end{tabular}
\end{center}

% \begin{itemize}
%     \item OFDMでは一般的にサブキャリア数$N \geq 50$
%     \item 最尤推定・固有モード伝送は$N$が大きくなればなるほど\TextCB{問題点がより厳しくなってくる}
%     \item 確率推論は\TextCA{$N$が大きいほうが有利}
% \end{itemize}
\vspace{0.2em}
\singlecolumn[0.07]{
\begin{infobox}{}
% \begin{itemize}
%     \item 基本的にはMIMO伝送の復調方式と同等な\\アルゴリズムによって復調可能
%     \item それぞれ一長一短があるのが現状
% \end{itemize}
\begin{itemize}
\item 最尤推定・固有モード伝送は$N$が大きくなればなるほど\\\TextCB{問題点がより厳しくなってくる}
% \item OFDMでは一般的にサブキャリア数$N \geq 50$
\item 確率推論は\TextCA{$N$が大きいほうが有利}
\end{itemize}
\end{infobox}
}
\end{frame}



\begin{frame}[t]{2022年10月RCS研究会での発表の概要}
\vspace{-0.7em}
\singlecolumn[0.11]{
% \vspace{-0.5em}
\begin{infobox}{目的}
サブキャリア数$N$の大きなSEFDMにおいて
\begin{itemize}
\item 確率推論手法の\TextCA{期待値伝播法（EP）}の性能評価
\item \TextCA{ハール・プリコーディング}による復調性能改善
\end{itemize}
\end{infobox}


\begin{infobox}{結果のまとめ}
\begin{itemize}
\item $N=160$, QPSK-SEFDMにおいてOFDMと\\比較して62.5\%まで帯域を削減可能
\end{itemize}
\end{infobox}
}
\end{frame}



% \begin{frame}[t]{SEFDMの変調方法の詳細}
%     \vspace{-0.5em}
%     \begin{center}
%         \small
%         \scalebox{0.8}{
%         \input{figs/sefdm_system.tex}
%         }
%     \end{center}

%     \vspace{-0.5em}


%     圧縮率$\alpha = N/K$のSEFDMの変調方法
%     \begin{enumerate}
%         \item $N$個のサブキャリアにプリコーディングをする
%         \item $K$点のIDFTに$N$点のサブキャリアを入力し\\\TextCA{先頭$N$サンプルだけ取得}（OFDM $\rightarrow$ SEFDM）
%         \item サイクリックプレフィックス（CP）を付与
%     \end{enumerate}
% \end{frame}


% \begin{frame}[t,fragile]{IDFTを利用したSEFDMの変調の概要}
%     \vspace{-0.5em}

%     \TextCA{ゼロパディングIDFT後のサブキャリア配置（OFDM）}
%     \begin{center}
%         \scalebox{2}{
%             \begin{tikzpicture}
%                 \pgfmathsetmacro\myPIE{3.141592}
%                 \pgfmathdeclarefunction{sinc}{1}{%
%                     \pgfmathparse{abs(#1)<0.001 ? int(1) : int(0)}%
%                     \ifnum\pgfmathresult>0 \pgfmathparse{1}\else\pgfmathparse{sin(#1 r * \myPIE)/(#1 * \myPIE) }\fi%
%                     }

%                 \draw (-1.5,0) -- (1.5,0);

%                 \foreach \s [count=\i] in {-2, -1, ..., 2}
%                 {
%                     \draw[domain=-3:3, samples=50] plot (\x*0.25+\s*0.25,{sinc(\x)*sinc(\x)});
%                     \draw [-stealth,CA](\s*0.25,0.95) -- (\s*0.25,0.05);
%                 }
%             \end{tikzpicture}
%         }
%     \end{center}

%     \TextCA{シンボル時間の短縮（サンプルの破棄）後（SEFDM）}
%     \begin{center}
%         \scalebox{2}{
%             \begin{tikzpicture}
%                 \pgfmathsetmacro\myPIE{3.141592}
%                 \pgfmathdeclarefunction{sinc}{1}{%
%                     \pgfmathparse{abs(#1)<0.001 ? int(1) : int(0)}%
%                     \ifnum\pgfmathresult>0 \pgfmathparse{1}\else\pgfmathparse{sin(#1 r * \myPIE)/(#1 * \myPIE) }\fi%
%                     }

%                 \draw (-1.5,0) -- (1.5,0);

%                 \foreach \s [count=\i] in {-2, -1, ..., 2}
%                 {
%                     \draw[domain=-3.5:3.5, samples=50] plot (\x*0.25+\s*0.25,{sinc(\x*0.5)*sinc(\x*0.5)});
%                     \draw [-stealth,CB](\s*0.25,0.95) -- (\s*0.25,0.50);
%                 }
%             \end{tikzpicture}
%         }
%     \end{center}
%     OFDMシンボルの末尾サンプルの破棄によりシンボル時間が短くなるため\TextCA{サブキャリア密度を維持したまま各サブキャリアの帯域幅が増加}$\rightarrow$ \TextCB{ICIの発生}
% \end{frame}


% \begin{frame}[t]{AWGN通信路でのSEFDMの行列表現}
%     \vspace{-2em}
%     \begin{align*}
%         \UnderlineAcap{$\mathmat{y}$}{受信} = \UnderlineAcap{$\mathmat{W}_{\alpha}$}{変調行列} \UnderlineAcap{$\mathmat{x}$}{送信} + \UnderlineAcap{$\mathmat{w}$}{雑音}
%     \end{align*}

%     \begin{itemize}
%         \item オーバーサンプリング率が1と仮定
%         \item 行列$\mathmat{W}_{\alpha}$はサイズ$N \times N$ \\
%             意味：$N$サブキャリアを時間軸$N$サンプルへ変換
%         \item 行列$\mathmat{W}_{\alpha}$の\TextCA{最大固有値と最小固有値の比（条件数）が1に近いほど$\mathmat{y}$から$\mathmat{x}$を推定しやすい}
%         \begin{itemize}
%             \item 圧縮率$\alpha=1$（OFDM）のとき条件数は1で\\
%             最も推定しやすい状況
%             \item \TextCB{圧縮率$\alpha$が小さいと条件数が大きく推定困難}
%         \end{itemize}
%     \end{itemize}
% \end{frame}



% \begin{frame}[t]{既存のSEFDMの復調方法}
%     \vspace{-2em}
%     \begin{center}
%         \renewcommand\arraystretch{2}
%         \tabcolsep = 0.1cm
%         \begin{tabular}{cccc}
%             基礎技術 & 復調性能 & 計算量 & 問題点 \\ \hline
%             最尤推定 & ○○ & ×× & 指数計算量 \\ \hline
%             固有モード伝送 & ○ & ○ & 小さな特異値の問題 \\ \hline
%             確率推論　& ○ & ○ & 小さな$N$で性能劣化 \\ \hline
%         \end{tabular}
%     \end{center}

%     \begin{itemize}
%         \item 基本的にはMIMO伝送の復調方式と同等な\\アルゴリズムによって復調可能
%         \item それぞれ一長一短があるのが現状
%     \end{itemize}
% \end{frame}




% \begin{frame}[t]{固有モード伝送（EVD）に基づく手法}
%     \vspace{-0.5em}
%     変調行列に固有モード伝送を適用
%     \begin{align*}
%         \mathmat{y} = \UnderlineA{$\mathmat{W}_{\alpha}$} \mathmat{x} + \mathmat{w} = \UnderlineA{$\mathmat{U} \mathmat{\Lambda} \mathmat{U}^{H}$} \mathmat{x} + \mathmat{w}
%     \end{align*}
%     送信信号$\mathmat{x}$に\UnderlineB{プリコーディング$\mathmat{U} \mathmat{\Lambda}^{-1}$}を適用して\\\TextCA{電力割当・固有モード伝送}をする
%     \begin{align*}
%         \mathmat{y} &= \UnderlineA{$\mathmat{U} \mathmat{\Lambda} \mathmat{U}^{H}$} (\UnderlineB{$\mathmat{U} \mathmat{\Lambda}^{-1}$} \mathmat{x}) + \mathmat{w} = \mathmat{U} \mathmat{x} + \mathmat{w}
%     \end{align*}
%     受信信号に行列$\mathmat{U}^{H}$を乗算するだけで推定可能
%     \begin{align*}
%         \mathmat{U}^{H} \mathmat{y} = \mathmat{x} + \mathmat{U}^{H} \mathmat{w}
%     \end{align*}
%     \vspace{-1.5em}

%     \begin{alertbox}{}
%         \centering
%         ランク落ちや固有値$\mathmat{\Lambda}$に小さな値がある場合に\\対応するモードでの伝送を諦める必要がある\\
%         \TextCB{帯域を狭くすればするほど伝送レードも低下}
%     \end{alertbox}
% \end{frame}


% \begin{frame}[t]{既存のSEFDMの復調方法}
%     \TextCA{最尤推定法（MLD）に基づく復調手法}
%     \begin{itemize}
%         \item 送信シンボルを網羅的に全探索する手法
%         \item サブキャリア数の\TextCB{指数関数で計算時間増加}
%         \item 計算量を削減した手法が存在 \\
%         \begin{itemize}
%             \item Sphere Decoding法
%             \item QR分解Mアルゴリズム（QRM-MLD）
%         \end{itemize}
%     \end{itemize}

%     \vspace{1em}

%     \TextCA{特異値分解（SVD）に基づく復調手法}
%     \begin{itemize}
%         \item 変調行列のSVD結果を利用し現実的な計算量で復調
%         \item 小さな特異値により\TextCB{伝送速度を犠牲にするか\\
%         誤り率を犠牲にする必要がある}
%     \end{itemize}
% \end{frame}




% \begin{frame}[t]{既存のSEFDMの復調方法（まとめ）}
%     \vspace{-2em}
%     \begin{center}
%         \renewcommand\arraystretch{2}
%         \tabcolsep = 0.1cm
%         \begin{tabular}{cccc}
%             基礎技術 & 復調性能 & 計算量 & 問題点 \\ \hline
%             最尤推定 & ○○ & ×× & \TextCB{指数計算量} \\ \hline
%             固有モード伝送 & ○ & ○ & \TextCB{小さな特異値の問題} \\ \hline
%             確率推論　& ○ & ○ & \textcolor{black!20}{小さな$N$で性能劣化} \\ \hline
%         \end{tabular}
%     \end{center}

%     % \begin{itemize}
%     %     \item 基本的にはMIMO伝送の復調方式と同等な\\アルゴリズムによって復調可能
%     %     \item それぞれ一長一短があるのが現状
%     % \end{itemize}
%     \begin{itemize}
%         \item 普通サブキャリア数$N \geq 64$がほとんど
%         \item 最尤推定・固有モード伝送は$N$が大きくなればなるほど\TextCB{問題点がより厳しくなってくる}
%         \item 確率推論は\TextCA{$N$が大きいほうが有利}
%     \end{itemize}
% \end{frame}




% \begin{frame}[t]{極端な例：圧縮率$\alpha=0$のとき}

% \end{frame}



\begin{frame}[t]{性能評価するSEFDMシステム（前回の発表より抜粋）}
\vspace{-0.5em}
\singlecolumn[0.1]{
\scalebox{1}{
\small
\input{figs/sefdm_eval_system.tex}
}


\vspace{1em}
\begin{itemize}
\item 送信サブキャリアに\TextCA{ハール直交行列}を作用
\item 受信側では\TextCA{期待値伝播法（EP）}で復調
\end{itemize}
}
\end{frame}



\begin{frame}[t]{AWGN，$N=160$，$\alpha=0.625$（前回の発表より抜粋）}
% \vspace{-0.5em}
\diffdoublecolumn{0.70}{0.35}{
\vspace{-2em}
\begin{center}
\includegraphics[page=3,width=0.95\textwidth]{plot_pdfs.pdf}
\end{center}
}{
{
\small
\begin{itemize}
\item \textcolor{C4}{QRM-MLD}：\\QR分解Mアルゴリズム\\最尤推定法
\item \textcolor{C3}{EMT}：固有モード伝送
\item \textcolor{C1}{GaBP}：ガウス信念伝播法
\item \textcolor{C2}{AMP}：\\近似的メッセージ伝播法
\item \textcolor{C0}{EP}：期待値伝播法
\end{itemize}
}
}



\vspace{-0.5em}
\centering{
\TextCA{EPは無圧縮OFDMの理論BER特性を達成}
}
% \end{infobox}
\end{frame}


\begin{frame}[t]{前回の発表の課題と本研究での解決方法}
\singlecolumn[0.07]{
\TextCB{復調処理（EP法）の計算量の問題}
\begin{itemize}
\item チャネル応答推定後に特異値分解$\mathcal{O}(N^3)$が必要
\item EP法の反復1回に$\mathcal{O}(N^2)$の計算コスト\\
（比較：OFDMは$\mathcal{O}(N\log N)$で復調可能）
\end{itemize}

\vspace{1em}
\begin{infobox}{本研究での解決策}
以下の計算が低コスト$\mathcal{O}(N\log N)$で実行可能な\\\TextCA{新しいSEFDMの変調行列$\mathmat{W}$の設計}
\begin{itemize}
\item 行列ベクトル積$\mathmat{W} \mathmat{x}$
\item 線形MMSEフィルタ$\mathmat{W}^{H} \brkts{ \sigma^2 \mathmat{I} + v \mathmat{W} \mathmat{W}^{H} }^{-1}$
\end{itemize}
\end{infobox}
}
\end{frame}



\begin{frame}[t]{本発表の概要}
\vspace{-0.7em}
\singlecolumn[0.09]{
% \vspace{-0.5em}
\begin{infobox}{目的}
サブキャリア数の大きなSEFDMにおいて
\begin{itemize}
\item 低計算量で期待値伝播法（EP）が適用可能な\\変調方法の考案：\TextCA{ランダムDFT拡散SEFDM}
\item 期待値伝播法（EP）適用時の誤り率の評価
\end{itemize}
\end{infobox}


\begin{infobox}{結果のまとめ}
\begin{itemize}
\item OFDMと同じチャネル推定等の技術が適用可能
\item FFTサイズ$N \geq 512$で良好な誤り率特性を達成
\item パスダイバーシチ効果を確認
\end{itemize}
\end{infobox}
}
\end{frame}



\begin{frame}[t]{SEFDM変調行列：特異値分布（$N=160$，$\alpha=10/16$）}
\diffdoublecolumn{0.62}{0.45}{
\vspace{-1em}
\includegraphics[width=1\textwidth,page=1]{plot_svd.pdf}
}{
圧縮率$\alpha$のSEFDM変調行列
\begin{align*}
&[\mathmat{W}_{\alpha}]_{m,n} = \exp\Brkts{-j\frac{2\pi}{N}mn {\color{red} \alpha} } \\
&\mathmat{W}_{\alpha} = \mathmat{U} \BackgroundB{$\displaystyle \mathmat{\Sigma}$} \mathmat{V}^{H} \\
&\BackgroundB{$\displaystyle \mathmat{\Sigma}$} = \mathrm{diag}\BRKTS{ \sigma_0, \sigma_1, \cdots, \sigma_{N-1}}\\
&\sigma_0 \geq \sigma_1 \geq \cdots \geq \sigma_{N-1}
\end{align*}

\begin{infobox}{}
$\alpha N$個の特異値のみ値を持ち\\$(1-\alpha)N$個の特異値が0
\end{infobox}
}
\end{frame}




\begin{frame}[t]{SEFDM変調行列：右ユニタリ行列（$N=160$，$\alpha=10/16$）}
\diffdoublecolumn{0.62}{0.45}{
\vspace{-1em}
\includegraphics[width=1\textwidth,page=2]{plot_svd.pdf}
}{
圧縮率$\alpha$のSEFDM変調行列
\begin{align*}
&[\mathmat{W}_{\alpha}]_{m,n} = \exp\Brkts{-j\frac{2\pi}{N}mn {\color{red} \alpha} } \\
&\mathmat{W}_{\alpha} = \mathmat{U} \sqrt{\mathmat{\Sigma}} \BackgroundB{$\displaystyle \sqrt{\mathmat{\Sigma}} \mathmat{V}^{H}$}\\
\end{align*}

\begin{infobox}{}
\begin{itemize}
\item 入力の$N$次元ベクトルを\\
\TextCA{$\alpha N$次元へ帯域圧縮}
\item サブキャリア間干渉
\end{itemize}
\end{infobox}
}
\end{frame}




\begin{frame}[t]{SEFDM変調行列：左ユニタリ行列（$N=160$，$\alpha=10/16$）}
\diffdoublecolumn{0.62}{0.45}{
\vspace{-1em}
\includegraphics[width=1\textwidth,page=3]{plot_svd.pdf}
}{
圧縮率$\alpha$のSEFDM変調行列
\begin{align*}
&[\mathmat{W}_{\alpha}]_{m,n} = \exp\Brkts{-j\frac{2\pi}{N}mn {\color{red} \alpha} } \\
&\BackgroundB{$\UnderlineAcap{$\mathmat{W}$}{DFT} \mathmat{U} \sqrt{\mathmat{\Sigma}}$} \sqrt{\mathmat{\Sigma}} \mathmat{V}^{H}\\
\end{align*}

\vspace{-1em}
\begin{infobox}{}
\begin{itemize}
\item DFTした結果が単位行列の$\alpha N$部分行列 \\
$\rightarrow$ \TextCA{IDFT行列の部分行列}
\end{itemize}
\end{infobox}
}
\end{frame}




\begin{frame}[t]{SEFDM変調行列と提案変調方法}
\headuline{CA}{ハール・プリコーディングSEFDM変調行列}
\begin{align*}
\mathmat{W}_{\alpha} \UnderlineC{$\mathmat{Q}$} = \UnderlineA{$\displaystyle \mathmat{U} \sqrt{\mathmat{\Sigma}}$} \UnderlineB{$ \displaystyle\sqrt{\mathmat{\Sigma}} \mathmat{V}^{H}$} \UnderlineC{$\mathmat{Q}$}
\end{align*}
\vspace{-2em}
\begin{itemize}
\item \UnderlineC{ハール行列$\mathmat{Q}$}：ランダムにサブキャリアを拡散
\item $\mathmat{W}_{\alpha}$の\UnderlineB{入力側$ \displaystyle\sqrt{\mathmat{\Sigma}} \mathmat{V}^{H}$}：$\alpha$倍の帯域圧縮（$\alpha < 1$）
\item $\mathmat{W}_{\alpha}$の\UnderlineA{出力側$\displaystyle \mathmat{U} \sqrt{\mathmat{\Sigma}}$}：逆離散フーリエ変換行列の部分行列
\end{itemize}

\vspace{1em}
\headuline{CA}{提案変調方法：ランダムDFT拡散SEFDM}
\begin{align*}
\mathmat{W}'_{\alpha} = \UnderlineAcap{$\mathmat{W}^{H}$}{IDFT} \UnderlineBcap{$\mathmat{D}$}{帯域圧縮} \UnderlineCcap{$\mathmat{P}_{\pi} \mathmat{W}$}{ランダムDFT拡散}
\end{align*}
\end{frame}


\begin{frame}[t]{提案変調方法の詳細}
\singlecolumn[0.05]{
\vspace{-2em}
\begin{align*}
\mathmat{W}'_{\alpha} = \UnderlineA{$\mathmat{W}^{H}$} \UnderlineB{$\mathmat{D}$} \UnderlineC{$\mathmat{P}_{\pi} \mathmat{W}$}
\end{align*}

\UnderlineC{\textbf{ランダムDFT拡散}$\mathmat{P}_{\pi} \mathmat{W}$}\headuline{CC}{}\\
\begin{itemize}
\item DFT行列の行をランダムに入れ替えた\textbf{ユニタリ行列}
\end{itemize}

\UnderlineB{\textbf{帯域圧縮}$\mathmat{D}$}\headuline{CB}{}\\
\begin{itemize}
\item $\alpha N$個だけ値を持つ\textbf{対角行列}
\end{itemize}

\UnderlineA{\textbf{IDFT行列}$\mathmat{W}^{H}$}\headuline{CA}{}\\
\begin{itemize}
\item OFDMと同じくIDFT行列で変調（\textbf{ユニタリ行列}）
\end{itemize}

\begin{infobox}{}
\begin{itemize}
\item<only@1> \TextCA{特異値分解と同じ構造}（ユニタリ行列で挟まれた対角行列）
\item<only@1> 三つの作用素すべて\TextCA{高々$\mathcal{O}(N \log N)$で計算可能}
\item<only@2> 周波数領域でランダム化・圧縮をしたDFT-s-OFDM
\item<only@2> \TextCA{OFDMと同じCP構造，周波数領域チャネル推定が可能}
\end{itemize}
\end{infobox}
}
\end{frame}



\begin{frame}[t]{提案変調方法のブロック図}
\vspace{-0.5em}
{
\small
\scalebox{0.95}{
\input{figs/rdft_s_sefdm_sys.tex}
}
}

\vspace{-1em}
\diffdoublecolumn{0.7}{0.35}{
\headuline{CA}{受信信号の周波数領域表現}
\begin{align*}
\mathmat{y} = \UnderlineC{$\mathmat{H} \mathmat{D} \mathmat{P}_{\pi} \mathmat{W}$} \mathmat{x} + \mathmat{w} = \UnderlineC{$\mathmat{A}$} \mathmat{x} + \mathmat{w}
\end{align*}
}{
\vspace{0.5em}
\rightnotebox{}{
$\mathmat{H}$：周波数応答対角行列\\
$\mathmat{w}$：AWGNベクトル\\
\UnderlineC{$\mathmat{A}$}：観測行列
}
}
\end{frame}




\begin{frame}[t]{期待値伝播法（EP）の概要}
\vspace{-0.5em}
\begin{center}
\small
\scalebox{1}{
\input{figs/ep_abstract.tex}
}
\end{center}


\begin{itemize}
\item 並列干渉除去と軟判定を交互に繰り返す手法
\item モジュールAは\TextCA{受信信号$\mathmat{y}$の情報を用いて送信シンボル$\mathmat{x}$を推定}
\begin{itemize}
\item 線形最小二乗誤差フィルタによる並列干渉除去
\end{itemize}
\item モジュールBは\TextCA{送信シンボル$\mathmat{x}$の変調の特徴を利用して$\mathmat{x}$を推定}
\begin{itemize}
\item AWGNとしたときの軟判定復調
\end{itemize}
\item 互いに推定値$\mathmat{x}$と分散$v$を交換して収束させる
\end{itemize}
\end{frame}


\begin{frame}[t]{期待値伝播法の処理}
\vspace{-0.5em}
\headuline{CA}{モジュールA：LMMSEフィルタでの並列干渉除去}
% \vspace{-1em}
\diffdoublecolumn{0.7}{0.5}{
% \scalebox{0.95}{
$\displaystyle \mathmat{x}_{\mathrm{A} \rightarrow \mathrm{B}}^t = \mathmat{x}_{\mathrm{B} \rightarrow \mathrm{A}}^t + \gamma_{v^t_{\mathrm{B} \rightarrow \mathrm{A}}} \mathmat{A}^H {\color{CB} \mathmat{\Xi}^{-1}_{v^t_{\mathrm{B} \rightarrow \mathrm{A}}} } \brkts{ \mathmat{y} - \mathmat{A} \mathmat{x}^t_{\mathrm{B} \rightarrow \mathrm{A}} }$ \\
$\displaystyle v_{\mathrm{A} \rightarrow \mathrm{B}}^{t} = \gamma_{v^t_{\mathrm{B} \rightarrow \mathrm{A}}} - v^t_{\mathrm{B} \rightarrow \mathrm{A}}$
% }
}{
% \scalebox{0.95}{
$\displaystyle \mathmat{\Xi}_{v} = \sigma^2 \mathmat{I}_M + v \mathmat{A}\mathmat{A}^H$\\
$\displaystyle \frac{1}{\gamma_v} = \frac{1}{N} \mathrm{Tr}\brkts{ {\color{CB} \mathmat{\Xi}_{v}^{-1}} \mathmat{A} \mathmat{A}^H }$
% }
}



\headuline{CA}{モジュールB：軟判定復調}
\diffdoublecolumn{0.5}{0.5}{
% \begin{align*}
$\displaystyle \mathmat{x}_{\mathrm{B}}^{t+1} = \mathbb{E}\BRKTS{\mathmat{x} | \mathmat{x}_{\mathrm{A} \rightarrow \mathrm{B}}^t }$\\
$\displaystyle \mathmat{x}_{\mathrm{B} \rightarrow \mathrm{A}}^{t+1} = v_{\mathrm{B} \rightarrow \mathrm{A}}^{t+1} \brkts{ \frac{\mathmat{x}_{\mathrm{B}}^{t+1}}{v_{\mathrm{B}}^{t+1}} - \frac{\mathmat{x}_{\mathrm{A} \rightarrow \mathrm{B}}^{t}}{v_{\mathrm{A} \rightarrow \mathrm{B}}^{t}}  }$
% \end{align*}
}{
% \begin{align*}
$\displaystyle v_\mathrm{B}^{t+1} = \frac{1}{N} \mathbb{E}\BRKTS{ \| \mathmat{x} - \mathmat{x}_\mathrm{B}^{t+1} \|^2  | \mathmat{x}_{\mathrm{A} \rightarrow \mathrm{B}}^t }$\\
$\displaystyle \frac{1}{v_{\mathrm{B} \rightarrow \mathrm{A}}^{t+1}} = \frac{1}{v_{\mathrm{B}}^{t+1}} - \frac{1}{v_{\mathrm{A} \rightarrow \mathrm{B}}^{t}}$
% \end{align*}
}



\begin{alertbox}{}
\begin{itemize}
\item モジュールAでの処理に\textbf{\TextCB{逆行列}}が必要，高計算量
\item 観測行列$\mathmat{A}$を事前に特異値分解$\mathcal{O}(N^3)$する方法などが考案
\end{itemize}
\end{alertbox}
\end{frame}

\begin{frame}[t]{提案変調方法における期待値伝播法}
\vspace{-0.5em}

\headuline{CA}{モジュールA：LMMSEフィルタでの並列干渉除去}
\diffdoublecolumn{0.7}{0.5}{
% \scalebox{0.95}{
$\displaystyle \mathmat{x}_{\mathrm{A} \rightarrow \mathrm{B}}^t = \mathmat{x}_{\mathrm{B} \rightarrow \mathrm{A}}^t + \gamma_{v^t_{\mathrm{B} \rightarrow \mathrm{A}}} \mathmat{A}^H {\color{CB} \mathmat{\Xi}^{-1}_{v^t_{\mathrm{B} \rightarrow \mathrm{A}}} } \brkts{ \mathmat{y} - \mathmat{A} \mathmat{x}^t_{\mathrm{B} \rightarrow \mathrm{A}} }$ \\
$\displaystyle v_{\mathrm{A} \rightarrow \mathrm{B}}^{t} = \gamma_{v^t_{\mathrm{B} \rightarrow \mathrm{A}}} - v^t_{\mathrm{B} \rightarrow \mathrm{A}}$
% }
}{
% \scalebox{0.95}{
$\displaystyle \mathmat{\Xi}_{v} = \sigma^2 \mathmat{I}_M + v \mathmat{A}\mathmat{A}^H$\\
$\displaystyle \frac{1}{\gamma_v} = \frac{1}{N} \mathrm{Tr}\brkts{ {\color{CB} \mathmat{\Xi}_{v}^{-1}} \mathmat{A} \mathmat{A}^H }$
% }
}


\begin{infobox}{}
\centering
提案変調方法における観測行列$\mathmat{A}$は\TextCA{\textbf{自然と特異値分解済み}}\\
になるように変調行列を設計（左ユニタリ行列は単位行列）
\vspace{-0.5em}
\begin{align*}
\mathmat{A} = \UnderlineBcap{$\mathmat{I}_N$}{左ユニタリ行列} \UnderlineAcap{$\mathmat{H} \mathmat{D}$}{特異値} \UnderlineBcap{$\mathmat{P}_{\pi} \mathmat{W}$}{右ユニタリ行列}
\end{align*}

\vspace{-0.5em}
\begin{itemize}
\item モジュールAで必要な逆行列${\color{CB} \mathmat{\Xi}_{v}^{-1}}$の計算は\TextCA{対角行列の逆行列になる}
\item 反復1回あたりDFT/IDFTがそれぞれ1回ずつ：\TextCA{$\mathcal{O}(N \log N)$}
\end{itemize}
\end{infobox}
\end{frame}


% \begin{frame}[t,noframenumbering]{EPのモジュールBでの処理}
%     \TextCA{モジュールB：軟判定復調}
%     \begin{align*}
%         \mathmat{x}_{\mathrm{B} \rightarrow \mathrm{A}}^{t+1} &= v_{\mathrm{B} \rightarrow \mathrm{A}}^{t+1} \brkts{ \frac{\mathmat{x}_{\mathrm{B}}^{t+1}}{v_{\mathrm{B}}^{t+1}} - \frac{\mathmat{x}_{\mathrm{A} \rightarrow \mathrm{B}}^{t}}{v_{\mathrm{A} \rightarrow \mathrm{B}}^{t}}  }, \\
%         \frac{1}{v_{\mathrm{B} \rightarrow \mathrm{A}}^{t+1}} &= \frac{1}{v_{\mathrm{B}}^{t+1}} - \frac{1}{v_{\mathrm{A} \rightarrow \mathrm{B}}^{t}}, \\
%         \mathmat{x}_{\mathrm{B}}^{t+1} &= \mathbb{E}\BRKTS{\mathmat{x} | \mathmat{x}_{\mathrm{A} \rightarrow \mathrm{B}}^t },\\
%         v_\mathrm{B}^{t+1} &= \frac{1}{N} \mathbb{E}\BRKTS{ \| \mathmat{x} - \mathmat{x}_\mathrm{B}^{t+1} \|^2  | \mathmat{x}_{\mathrm{A} \rightarrow \mathrm{B}}^t }
%     \end{align*}
% \end{frame}

% \begin{frame}[t]{提案変調方法における期待値伝播法の低計算量化}

% \end{frame}

\begin{frame}[t]{性能評価シミュレーション諸元}
\begin{center}
\begin{tabular}{c|c}
\hline
% Parameter & Value \\ \hline
サブキャリア変調 & QPSK \\ \hline
FFTサイズ & 64, 256, 512, 1024, 2048 \\ \hline
圧縮率 & 0.625 = 10/16 \\ \hline
サブキャリア数 & 40, 160, 320, 640, 1280 \\ \hline
通信路 & AWGN，周波数選択性レイリーフェージング \\ \hline
\end{tabular}
\end{center}

\vspace{1em}
\headuline{CA}{独立$L$波周波数選択性レイリーフェージング}
\begin{align*}
h[n] = \begin{cases}
\mathcal{CN}(0, 1/L) & n < L \\
0 & \text{otherwise}
\end{cases}
\end{align*}
\end{frame}


\begin{frame}[t]{AWGN通信路における性能評価}
\diffdoublecolumn{0.67}{0.38}{
\vspace{-1.4em}
\includegraphics[page=1,width=\textwidth]{plot_ber.pdf}
}{
圧縮率$\alpha = 0.625$ \\
繰返し回数$I = 20$ \\

\begin{infobox}{}
\begin{itemize}
\item $N=64$のとき\\$10^{-2}$付近でフロア
\item $N=256$のとき\\
$10^{-5}$付近でフロア
\item $N\geq 512$のとき\\
\TextCA{フロアが見られない\\
良好特性}
\end{itemize}
\end{infobox}
}
\end{frame}



\begin{frame}[t]{独立$L$波レイリーフェージング通信路における性能評価}
\diffdoublecolumn{0.67}{0.38}{
\vspace{-1.4em}
\includegraphics[page=2,width=\textwidth]{plot_ber.pdf}
}{
% 圧縮率$\alpha = 0.625$ \\
FFTサイズ$N=512$\\
繰返し回数$I = 20$ \\

\begin{infobox}{}
\begin{itemize}
\item 独立波の数を増やすと
性能改善
\item 周波数領域でのランダムDFT拡散による\TextCA{パスダイバーシチ効果}
\end{itemize}
\end{infobox}
}
\end{frame}



\begin{frame}[t]{EP法の繰り返し回数による性能評価}
\diffdoublecolumn{0.67}{0.38}{
\vspace{-1.4em}
\includegraphics[page=4,width=\textwidth]{plot_ber.pdf}
}{
% 圧縮率$\alpha = 0.625$ \\
FFTサイズ$N=512$\\
独立波数$L=8$ \\

\begin{infobox}{}
\begin{itemize}
\item 12回以上繰り返すとほとんど収束\\
\item OFDMの24倍程度の計算量
\end{itemize}
\end{infobox}
}
\end{frame}



\begin{frame}[t]{本発表のまとめ}
\vspace{-0.7em}
\singlecolumn[0.07]{
% \vspace{-0.5em}
\begin{infobox}{目的}
\centering
サブキャリア数の大きなSEFDMのための\\
新しい低計算量な変調・復調方法の提案
\end{infobox}


\begin{infobox}{結果のまとめ}
\begin{itemize}
\item $N \geq 512$で良好な誤り率特性を達成
\item パスダイバーシチ効果を確認
\end{itemize}
\end{infobox}
\begin{alertbox}{今後の課題}
チャネル推定誤差や高いPAPRによる非線形性の影響の調査
\end{alertbox}
}
\end{frame}





% \begin{frame}[t,noframenumbering]{SEFDMのAWGN通信路システムモデル}

%     \diffdoublecolumn{0.6}{0.45}
%     {
%         \vspace{-1em}
%         \begin{center}
%             \scalebox{0.9}
%             {
%                 \input{figs/sefdm_awgn.tex}
%             }
%         \end{center}
%     }
%     {
%         \begin{align*}
%             &\mathmat{y} = \mathmat{W}_{\alpha} \mathmat{x} + \mathmat{w} \\
%             &[\mathmat{W}_{\alpha}]_{m,n} = \exp\Brkts{-j\frac{2\pi}{N}mn {\color{red} \alpha} }
%         \end{align*}
%     }

%     \vspace{2em}


%     \diffdoublecolumn{0.78}{0.27}
%     {
%         \begin{itemize}
%             \item $\alpha = 1$でOFDMと等価（$\mathmat{W}_{\alpha=1}$がIDFT行列に一致）
%             \item $\alpha = 0$は全サブキャリアを同一のキャリアに重畳
%             \item \TextCA{$0 < \alpha < 1$で受信信号から送信信号を推定したい}
%         \end{itemize}
%     }
%     {
%         \rightnotebox{Note}{
%             \begin{align*}
%                 &\text{圧縮率 } \alpha = N/K < 1 \\
%                 &\text{サブキャリア数} N \\
%                 &\text{FFTサイズ} K
%             \end{align*}
%         }
%     }

% \end{frame}



% \begin{frame}[t,noframenumbering]{まとめ}
%     % \vspace{-0.5em}
%     \begin{infobox}{本研究の目的}
%         サブキャリア数$N$の大きなSEFDMにおいて
%         \begin{itemize}
%             \item 確率推論手法の\TextCA{期待値伝播法（EP）}の性能評価
%             \item \TextCA{ハール・プリコーディング}による復調性能改善
%         \end{itemize}
%     \end{infobox}

%     \begin{infobox}{結果のまとめ}
%         \begin{itemize}
%             \item $N=160$, QPSK-SEFDMにおいてOFDMと\\比較して62.5\%まで帯域を削減可能
%             \item プリコーディングなしではAWGNでも\\全く復調できない
%         \end{itemize}
%     \end{infobox}
% \end{frame}






% \begin{frame}[t,noframenumbering]{今後の課題}
%     \TextCB{復調処理（EP法）の計算量の問題}
%     \begin{itemize}
%         \item 通信路特性推定後に特異値分解（$\mathcal{O}(N^3)$）が必要
%         \item EP法の反復1回に$\mathcal{O}(N^2)$の計算コスト\\
%         （比較：OFDMは$\mathcal{O}(N\log N)$で復調可能）
%     \end{itemize}

%     \vspace{1em}

%     \TextCA{解決方法}\\
%     以下の計算が低コスト（$\mathcal{O}(N\log N)$）で実行可能な\\\TextCA{新しいSEFDMの変調行列$\mathmat{W}$の設計}
%     \begin{itemize}
%         \item 行列ベクトル積$\mathmat{W} \mathmat{x}$
%         \item 線形MMSEフィルタ$\mathmat{W}^{H} \brkts{ \sigma^2 \mathmat{I} + v \mathmat{W} \mathmat{W}^{H} }^{-1}$
%     \end{itemize}

%     % \begin{infobox}{}
%     %     \begin{itemize}
%     %         \item SEFDMと同等な周波数領域での圧縮を達成する新しい変調方法
%     %     \end{itemize}
%     % \end{infobox}
% \end{frame}



% \begin{frame}[t,noframenumbering]{まとめ}
%     % \vspace{-0.5em}
%     \begin{infobox}{本研究の目的}
%         サブキャリア数$N$の大きなSEFDMにおいて
%         \begin{itemize}
%             \item 確率推論手法の\TextCA{期待値伝播法（EP）}の性能評価
%             \item \TextCA{ハール・プリコーディング}による復調性能改善
%         \end{itemize}
%     \end{infobox}

%     \begin{infobox}{結果のまとめ}
%         \begin{itemize}
%             \item $N=160$, QPSK-SEFDMにおいてOFDMと\\比較して62.5\%まで帯域を削減可能
%             \item プリコーディングなしではAWGNでも\\全く復調できない
%         \end{itemize}
%     \end{infobox}
% \end{frame}

\begin{frame}[t,noframenumbering]{EP法とハール・プリコーディング}

\singlecolumn[0.07]{
\vspace{-0.5em}

% \begin{infobox}{ハール直交行列}
\headuline{CA}{ハール直交行列} \\
要素が複素正規分布する行列$\mathmat{X}$をQR分解したときの\\\TextCA{ユニタリ行列$\mathmat{Q}$はハール直交行列}
\begin{align*}
\mathmat{X} = \UnderlineC{$\mathmat{Q}$} \mathmat{R}
\end{align*}

\vspace{1em}
% \end{infobox}
% \vspace{1em}

\headuline{CA}{ハール・プリコーディング}
\begin{itemize}
\item 送信側でハール直交行列を作用
\item 大規模MIMOにおいて大幅にEP法の復調性能改善
\end{itemize}


\begin{infobox}{}
\centering
SEFDMの復調においても性能改善されるのか調査
\end{infobox}
}
\end{frame}



\begin{frame}[t,noframenumbering]{EPのモジュールAでの処理}
\TextCA{モジュールA：LMMSEフィルタでの並列干渉除去}
\begin{align*}
\mathmat{x}_{\mathrm{A} \rightarrow \mathrm{B}}^t &= \mathmat{x}_{\mathrm{B} \rightarrow \mathrm{A}}^t + \gamma_{v^t_{\mathrm{B} \rightarrow \mathrm{A}}} \mathmat{A}^H \mathmat{\Xi}^{-1}_{v^t_{\mathrm{B} \rightarrow \mathrm{A}}} \brkts{ \mathmat{y} - \mathmat{A} \mathmat{x}^t_{\mathrm{B} \rightarrow \mathrm{A}} }, \\
v_{\mathrm{A} \rightarrow \mathrm{B}}^{t} &= \gamma_{v^t_{\mathrm{B} \rightarrow \mathrm{A}}} - v^t_{\mathrm{B} \rightarrow \mathrm{A}}, \\
\mathmat{\Xi}_{v} &= \sigma^2 \mathmat{I}_M + v \mathmat{A}\mathmat{A}^H,\\
\frac{1}{\gamma_v} &= \frac{1}{N} \mathrm{Tr}\brkts{ \mathmat{\Xi}_{v}^{-1} \mathmat{A} \mathmat{A}^H }
\end{align*}
\end{frame}
\begin{frame}[t,noframenumbering]{EPのモジュールBでの処理}
\TextCA{モジュールB：軟判定復調}
\begin{align*}
\mathmat{x}_{\mathrm{B} \rightarrow \mathrm{A}}^{t+1} &= v_{\mathrm{B} \rightarrow \mathrm{A}}^{t+1} \brkts{ \frac{\mathmat{x}_{\mathrm{B}}^{t+1}}{v_{\mathrm{B}}^{t+1}} - \frac{\mathmat{x}_{\mathrm{A} \rightarrow \mathrm{B}}^{t}}{v_{\mathrm{A} \rightarrow \mathrm{B}}^{t}}  }, \\
\frac{1}{v_{\mathrm{B} \rightarrow \mathrm{A}}^{t+1}} &= \frac{1}{v_{\mathrm{B}}^{t+1}} - \frac{1}{v_{\mathrm{A} \rightarrow \mathrm{B}}^{t}}, \\
\mathmat{x}_{\mathrm{B}}^{t+1} &= \mathbb{E}\BRKTS{\mathmat{x} | \mathmat{x}_{\mathrm{A} \rightarrow \mathrm{B}}^t },\\
v_\mathrm{B}^{t+1} &= \frac{1}{N} \mathbb{E}\BRKTS{ \| \mathmat{x} - \mathmat{x}_\mathrm{B}^{t+1} \|^2  | \mathmat{x}_{\mathrm{A} \rightarrow \mathrm{B}}^t }
\end{align*}
\end{frame}

% \begin{frame}[t,noframenumbering]{周波数選択性レイリーフェージング}
%     \vspace{-2em}
%     \begin{center}
%     \scalebox{0.5}
%     {
%         \input{../plot/ber_256_160_Rayleigh.pgf}
%     }
%     \end{center}


%     \begin{itemize}
%         \item<only@1> サブキャリア数$N=160$，圧縮率$\alpha=0.625$
%         \item<only@1> EP法のみ復調可能
%         \item<only@2> 純OFDMの理論より良い復調特性\\
%         プリコーディングにより周波数方向に拡散
%     \end{itemize}
% \end{frame}



\begin{frame}[t,noframenumbering]{最尤推定法（MLD）に基づく復調手法}
送信シンボルを網羅的に全探索する手法
\begin{itemize}
\item すべての候補で以下のコスト関数を計算
\begin{align*}
\mathcal{J}(\mathmat{x}) = \| \mathmat{y} - \mathmat{W}_{\alpha} \mathmat{x} \|^2
\end{align*}
\item コスト関数が最小の候補を送信信号と推定
\item サブキャリア数$N$の\TextCB{指数関数で計算時間増加} \\
例１：QPSKの$N$サブキャリアで$4^N$個の候補\\
例２：BPSKで$N = 2$なら4個の候補
\begin{align*}
\mathmat{x} \in \Brkts{ \begin{bmatrix}+1\\+1\end{bmatrix}, \begin{bmatrix}-1\\+1\end{bmatrix}, \begin{bmatrix}+1\\-1\end{bmatrix}, \begin{bmatrix}-1\\-1\end{bmatrix} }
\end{align*}
\end{itemize}
\end{frame}


\begin{frame}[t,fragile,noframenumbering]{固有モード伝送（EMT）法}
\only<1>{
\vspace{-0.5em}
\scalebox{0.85}{
\small
\input{figs/sefdm_evd.tex}
}

\begin{itemize}
\item 変調行列の特異値分解（SVD）$\mathmat{W}_{\alpha} = \mathmat{U} \mathmat{S} \mathmat{V}^{H}$
\item 特異ベクトル$\mathmat{U}, \mathmat{V}$でプリ/ポストコーディング
\item 特異値の逆数$\mathmat{S}^{-1}$で最適電力割当
\end{itemize}
}

\only<2>{
\begin{infobox}{EMT法の利点}
\begin{itemize}
\item 送受信機での処理が簡単
\item チャネル推定が完全なら誤り率特性も最適
\end{itemize}
\end{infobox}
\vspace{1em}
\begin{alertbox}{EMT法の欠点}
\begin{itemize}
\item ランク落ちや特異値に小さな値がある場合に\\対応するモードでの伝送を諦める必要がある\\
\item \TextCB{帯域を狭くすればするほど伝送レードも低下}
\end{itemize}
\end{alertbox}
}
\end{frame}


\begin{frame}[t,noframenumbering]{変調行列の特異値分布（$N=256$）}
\vspace{-2em}
% \begin{minipage}[t]{0.49\textwidth}
\centering
\scalebox{0.50}{
\input{../plot/eigen_values_256.pgf}
}



\begin{itemize}
\item 帯域を圧縮率するほど小さな特異値の割合が増加
\item 計算誤差によるEVMの増加$\rightarrow$ ビット誤り率の増大
\end{itemize}
\end{frame}


\begin{frame}[t,noframenumbering]{EMT法のエラーベクトル振幅（EVM）}
\vspace{-2em}
% \begin{minipage}[t]{0.49\textwidth}
\centering
\scalebox{0.50}{
\input{../plot/evd_evm.pgf}
}

% \begin{itemize}
%     \item 帯域を圧縮率するほど小さな固有値の割合が増加
%     \item 計算誤差によるEVMの増加$\rightarrow$ ビット誤り率の増大
% \end{itemize}
\vspace{-0.5em}
\begin{itemize}
\item サブキャリア数が大きいと少しの圧縮でEVM劣化
\item \TextCB{$N=256$ではほとんど圧縮不可}
\end{itemize}
\end{frame}


\begin{frame}[t,noframenumbering]{既存のSEFDMの復調方法（まとめ）}
\vspace{-2em}
\begin{center}
\renewcommand\arraystretch{2}
\tabcolsep = 0.1cm
\begin{tabular}{cccc}
基礎技術 & 復調性能 & 計算量 & 問題点 \\ \hline
最尤推定 & ○○ & ×× & \TextCB{指数計算量} \\ \hline
固有モード伝送 & ○ & ○ & \TextCB{小さな特異値の問題} \\ \hline
確率推論　& ○ & ○ & \textcolor{black!20}{小さな$N$で性能劣化} \\ \hline
\end{tabular}
\end{center}

% \begin{itemize}
%     \item 基本的にはMIMO伝送の復調方式と同等な\\アルゴリズムによって復調可能
%     \item それぞれ一長一短があるのが現状
% \end{itemize}
\begin{itemize}
\item 一般的なOFDMのサブキャリア数$N \geq 50$
\item 最尤推定・固有モード伝送は$N$が大きくなればなるほど\TextCB{問題点がより厳しくなってくる}
\item 確率推論は\TextCA{$N$が大きいほうが有利}
\end{itemize}
\end{frame}



\begin{frame}[t,noframenumbering]{確率推論に基づく手法}
% \begin{infobox}{AWGN線形通信路モデル}
%     ある線形システム$\mathmat{A}$とAWGN通信路$\mathmat{w}$を通って受信された受信ベクトル$\mathmat{y}$から送信信号ベクトル$\mathmat{x} \sim p(\mathmat{x})$を推定したい
%     \begin{align*}
%         \mathmat{y} = \mathmat{z} + \mathmat{w};\;\;\;\; \mathmat{z} = \mathmat{A} \mathmat{x}
%     \end{align*}
% \end{infobox}

\begin{infobox}{確率推論}
送信信号$\mathmat{x}$の事前確率$p(\mathmat{x})$や通信路の確率モデル$\mathmat{p}(\mathmat{y}|\mathmat{z} = \mathmat{A}\mathmat{x})$を利用して送信信号$\mathmat{x}$を推定する
\end{infobox}

確率モデル上でのメッセージ伝播アルゴリズム
\begin{itemize}
\item 並列干渉除去と信号推定を繰り返し動作
\item バリエーションが存在（GaBP, AMP, EP）
\item 大システム極限の仮定\\
\TextCA{サブキャリア数$N$が大きいときに性能発揮}
% \begin{itemize}
%     \item ガウス信念伝播法（GaBP）
%     \item 近似的メッセージ伝播法（AMP）
%     \item 期待値伝播法（EP）
% \end{itemize}
\end{itemize}


% \vspace{1em}

% \begin{alertbox}{}
%     \centering
%     大システムを仮定しているのでサブキャリア数$N$が\\小さいときに著しく性能劣化
% \end{alertbox}
\end{frame}




\begin{frame}[t,noframenumbering]{性能評価：諸元}
\vspace{-0.5em}
% {
% \centering
\begin{center}
\begin{tabular}{c|c}
\hline
% Parameter & Value \\ \hline
サブキャリア変調 & QPSK \\ \hline
FFTサイズ & 16, 64, 256 \\ \hline
圧縮率 & 0.625 \\ \hline
サブキャリア数 & 10, 40, 160 \\ \hline
\end{tabular}
\end{center}
% }

{
以下の手法でビット誤り率を比較
\begin{itemize}
\item 最尤推定法（MLD）
\item QR分解MアルゴリズムMLD法（QRM-MLD）
\item 固有モード伝送法（EMT）
\item 確率推論系
\begin{itemize}
\item ガウス信念伝播法（GaBP）
\item 近似的メッセージ伝播法（AMP）
\item 期待値伝播法（EP）
\end{itemize}
\end{itemize}
}
\end{frame}


\begin{frame}[t,noframenumbering]{AWGN，$N=10$，$\alpha=0.625$　}
\vspace{-1em}
\begin{center}
% \scalebox{0.58}{
% \input{../plot/ber_16_10_AWGN_forslide.pgf}
% }
\includegraphics[page=1,width=0.74\textwidth]{plot_pdfs.pdf}
\end{center}

\vspace{-1.5em}
\begin{itemize}
\item<only@1> 小規模SEFDMで\textcolor{C3}{固有モード伝送（EMT）}は\\無圧縮のOFDMと同じBER特性を達成
\item<only@1> \textcolor{C5}{最尤推定（MLD）系}はOFDMより少し劣化
\item<only@2> 確率推論系は\TextCB{全く復調できない}\\
$\rightarrow$ 大システム極限の仮定が成立しないため
\end{itemize}
\end{frame}


\begin{frame}[t,noframenumbering]{AWGN，$N=40$，$\alpha=0.625$　}
\vspace{-1em}
\begin{center}
\includegraphics[page=2,width=0.74\textwidth]{plot_pdfs.pdf}
\end{center}

\vspace{-1.5em}
\begin{itemize}
\item<only@1> 中規模SEFDMではどの手法も復調困難
\item<only@1> MLD系は計算量・探索空間の問題によりBER増加
\item<only@1> EMT法はEVMの劣化によるBER増加
\item<only@2> サブキャリア数$N=40$程度では確率推論系も\\まだうまく動かない
\end{itemize}
\end{frame}


\begin{frame}[t,noframenumbering]{ハール・プリコーディング有無の比較}
\vspace{-1em}
\begin{center}
\includegraphics[page=4,width=0.74\textwidth]{plot_pdfs.pdf}
\end{center}

\vspace{-1.5em}
\begin{itemize}
\item 大規模SEFDM（$N=160$, $\alpha=0.625$）で送信側のハール・プリコーディングの有無でBERを比較
\item \TextCA{プリコーディングにより復調性能が大幅に改善}
\end{itemize}
\end{frame}
